\chapter{Numeros naturales e Induccion}

\section{Sucesiones}

\begin{definicion}[title= Definicion de sucesion]
    Una sucesion es una lista ordenada de elementos que siguen una regla o "patron".
\end{definicion}

\begin{definicion}[title= Definicion de serie]
    Una serie es la suma de todos los elementos de una sucesion infinita.
\end{definicion}

\begin{definicion}[title= Definicion de serie geometrica]
    Es la serie que toma las $n+1$ primeras potencias de un numero ($q$). \\
    Si aplicamos la logica de la suma de Gauss nos queda algo de la forma:
\[
    \sum_{k=0}^{n} q^k = \frac{q^{n+1}-1}{q-1}, \quad \text{si } q\neq 1
\]
\[
    \sum_{k=0}^{n} q^k = n+1, \quad \text{si } q=1
\]
\end{definicion}

\begin{definicion}[title= Definicion de serie telescopica]
    Es una serie donde cada termino se puede escribir como una diferencia $a_k - a_{k+1}$, de modo que casi todos los terminos se cancelan entre si:
\[
    \sum_{k=1}^{n} (a_k - a_{k+1}) = a_1 - a_{n+1}
\]
\end{definicion}

\begin{definicion}[title= Definicion suma de Gauss]
    Es la suma de los primeros $n$ numeros naturales. Se obtiene sumando el primero con el ultimo, el segundo con el anteultimo, etc.:
\[
    \sum_{k=1}^{n} k = \frac{n(n+1)}{2}
\]
\end{definicion}

\begin{definicion}[title= Definicion sucesion de Lucas]
    Es la sucesion definida por recurrencia con la misma regla que Fibonacci pero con otros valores iniciales:
\[
    L_0 = 2, \quad L_1 = 1, \quad L_{n} = L_{n-1} + L_{n-2} \ \text{ para } n\geq 2
\]
    Sus primeros terminos son $2, 1, 3, 4, 7, 11, 18, \dots$
\end{definicion}

\section{Induccion}

\begin{definicion}[title= Principio de induccion 1]
    Sea $P(n)$ una proposicion sobre $n\in\mathbb{N}$. Si se cumple que:
    \begin{itemize}
        \item (Caso base) $P(n_0)$ es verdadera.
        \item (Paso inductivo) Para todo $k\geq n_0$: $P(k) \Rightarrow P(k+1)$.
    \end{itemize}
    Entonces $P(n)$ es verdadera para todo $n\geq n_0$.
\end{definicion}

\begin{definicion}[title= Principio de induccion 2]
    Sea $P(n)$ una proposicion sobre $n\in\mathbb{N}$. Si se cumple que:
    \begin{itemize}
        \item (Casos base) $P(n_0)$ y $P(n_0+1)$ son verdaderas.
        \item (Paso inductivo) Para todo $k\geq n_0$: $P(k) \wedge P(k+1) \Rightarrow P(k+2)$.
    \end{itemize}
    Entonces $P(n)$ es verdadera para todo $n\geq n_0$. Es util para recurrencias de dos terminos (Fibonacci, Lucas).
\end{definicion}

\begin{definicion}[title= Induccion fuerte]
    Sea $P(n)$ una proposicion sobre $n\in\mathbb{N}$. Si se cumple que:
    \begin{itemize}
        \item (Caso base) $P(n_0)$ es verdadera.
        \item (Paso inductivo) Para todo $k\geq n_0$: si $P(j)$ es verdadera para todo $n_0\leq j\leq k$, entonces $P(k+1)$ es verdadera.
    \end{itemize}
    Entonces $P(n)$ es verdadera para todo $n\geq n_0$.
\end{definicion}